% =====================================================================
% Agentic FedGrid -- LaTeX-ready section text (IEEE conference style).
% Tables come from \begin{table}[t]\centering
\caption{Frozen backbone before federation: official UNSW-NB15 test set (82{,}332 flows) and ToN\_IoT held-out split (31{,}656 flows). v1: one token for all continuous features, 10 epochs, single run. v2: one token per feature, 30 epochs, one-cycle LR, mean $\pm$ sd over 3 seeds. Trained and evaluated on EC2.}
\label{tab:backbone}
\begin{tabular}{lcc}\toprule
Metric & v1 & v2 \\ \midrule
UNSW accuracy (\%) & 87.19 & 87.73 $\pm$ 0.35 \\
UNSW macro-F1 & 0.871 & 0.876 $\pm$ 0.004 \\
UNSW FPR (\%) & 14.31 & 12.94 $\pm$ 1.13 \\
Fuzzers recall (\%) & 41.06 & 36.81 $\pm$ 0.81 \\
ToN\_IoT accuracy (\%) & 98.84 & 99.28 $\pm$ 0.04 \\
ToN\_IoT FPR (\%) & 1.93 & 1.64 $\pm$ 0.11 \\
Median latency (ms, t3.small) & 1.71 & 1.70 $\pm$ 0.02 \\
\bottomrule\end{tabular}\end{table}

\begin{table}[t]\centering
\caption{Global model after each FedAvg round, backbone v1 (mean $\pm$ sd over 4 seeds; real Lambda handler code on EC2 t3.small with DynamoDB/S3 mocked).}
\label{tab:fed_rounds_v1}
\begin{tabular}{ccccc}\toprule
Round & Fuzzers recall (\%) & Replay acc. (\%) & UNSW acc. (\%) & UNSW FPR (\%) \\ \midrule
0 & 41.06 $\pm$ 0.00 & 28.14 $\pm$ 0.00 & 87.19 $\pm$ 0.00 & 14.31 $\pm$ 0.00 \\
1 & 41.36 $\pm$ 0.87 & 28.39 $\pm$ 1.08 & 87.36 $\pm$ 0.25 & 14.05 $\pm$ 0.87 \\
2 & 45.36 $\pm$ 0.76 & 32.57 $\pm$ 0.92 & 86.82 $\pm$ 0.04 & 16.44 $\pm$ 0.23 \\
3 & 48.12 $\pm$ 2.13 & 35.90 $\pm$ 2.55 & 86.72 $\pm$ 0.08 & 17.20 $\pm$ 0.73 \\
4 & 54.50 $\pm$ 1.92 & 43.53 $\pm$ 1.98 & 86.46 $\pm$ 0.07 & 19.21 $\pm$ 0.61 \\
\bottomrule\end{tabular}\end{table}

\begin{table}[t]\centering
\caption{Global model after each FedAvg round, backbone v2 (mean $\pm$ sd over 3 seeds; real Lambda handler code on EC2 t3.small with DynamoDB/S3 mocked).}
\label{tab:fed_rounds_v2}
\begin{tabular}{ccccc}\toprule
Round & Fuzzers recall (\%) & Replay acc. (\%) & UNSW acc. (\%) & UNSW FPR (\%) \\ \midrule
0 & 36.81 $\pm$ 0.81 & 24.17 $\pm$ 1.02 & 87.73 $\pm$ 0.35 & 12.94 $\pm$ 1.13 \\
1 & 37.83 $\pm$ 1.26 & 25.63 $\pm$ 1.24 & 88.02 $\pm$ 1.07 & 12.66 $\pm$ 2.72 \\
2 & 41.80 $\pm$ 2.08 & 30.15 $\pm$ 2.54 & 87.53 $\pm$ 0.36 & 15.16 $\pm$ 0.48 \\
3 & 44.22 $\pm$ 5.97 & 32.55 $\pm$ 7.03 & 87.25 $\pm$ 0.03 & 16.46 $\pm$ 1.97 \\
4 & 47.39 $\pm$ 9.83 & 36.35 $\pm$ 11.74 & 87.30 $\pm$ 0.76 & 17.01 $\pm$ 4.69 \\
\bottomrule\end{tabular}\end{table}

\begin{table*}[t]\centering
\caption{Ablations after 4 federated rounds, global model on device 2 (mean $\pm$ sd over 3 seeds; EC2). Payload = raw float32 bytes uploaded per device per round. Adv.\ acc.\ = accuracy under FGSM, $\epsilon{=}0.05$, 5{,}000 test flows.}
\label{tab:ablation}
\begin{tabular}{lrcccccc}\toprule
Variant & Payload (B) & Fuzzers rec. (\%) & UNSW acc. (\%) & UNSW FPR (\%) & Macro-F1 & Adv. acc. (\%) & ToN acc. (\%) \\ \midrule
LoRA $r{=}1$ & 2,835 & 51.25 $\pm$ 1.71 & 86.67 $\pm$ 0.13 & 18.26 $\pm$ 0.69 & 0.864 $\pm$ 0.001 & 73.82 $\pm$ 1.64 & 98.90 $\pm$ 0.04 \\
LoRA $r{=}2$ & 5,403 & 52.42 $\pm$ 0.11 & 86.55 $\pm$ 0.14 & 18.66 $\pm$ 0.30 & 0.863 $\pm$ 0.001 & 73.50 $\pm$ 1.48 & 98.87 $\pm$ 0.01 \\
LoRA $r{=}4$ (default) & 10,515 & 50.48 $\pm$ 3.16 & 86.47 $\pm$ 0.04 & 18.28 $\pm$ 0.83 & 0.862 $\pm$ 0.001 & 74.49 $\pm$ 0.15 & 98.19 $\pm$ 1.09 \\
LoRA $r{=}8$ & 20,763 & 59.12 $\pm$ 6.45 & 86.00 $\pm$ 0.56 & 21.24 $\pm$ 2.66 & 0.857 $\pm$ 0.007 & 75.21 $\pm$ 1.07 & 97.72 $\pm$ 0.97 \\
Full-model FedAvg & 106,155 & 78.04 $\pm$ 3.34 & 82.66 $\pm$ 0.40 & 33.08 $\pm$ 1.80 & 0.817 $\pm$ 0.005 & 59.13 $\pm$ 13.48 & 95.95 $\pm$ 2.51 \\
w/o hard-example mining & 10,515 & 52.97 $\pm$ 2.58 & 86.62 $\pm$ 0.02 & 18.65 $\pm$ 0.70 & 0.864 $\pm$ 0.000 & 58.71 $\pm$ 12.80 & 98.81 $\pm$ 0.05 \\
w/o FGSM training & 10,515 & 50.26 $\pm$ 0.15 & 86.78 $\pm$ 0.03 & 17.67 $\pm$ 0.12 & 0.866 $\pm$ 0.000 & 52.37 $\pm$ 0.37 & 99.11 $\pm$ 0.01 \\
FedProx ($\mu{=}0.01$) & 10,515 & 52.42 $\pm$ 3.88 & 86.37 $\pm$ 0.10 & 18.98 $\pm$ 1.13 & 0.861 $\pm$ 0.001 & 74.56 $\pm$ 0.16 & 98.81 $\pm$ 0.01 \\
Median aggregation & 10,515 & 56.84 $\pm$ 3.00 & 86.31 $\pm$ 0.15 & 20.14 $\pm$ 1.02 & 0.860 $\pm$ 0.002 & 75.01 $\pm$ 0.89 & 98.21 $\pm$ 1.05 \\
+ adoption gate & 10,515 & 48.88 $\pm$ 4.95 & 86.47 $\pm$ 0.12 & 17.83 $\pm$ 1.43 & 0.863 $\pm$ 0.002 & 74.58 $\pm$ 0.34 & 98.18 $\pm$ 1.08 \\
Poisoned client, FedAvg & 10,515 & 76.20 $\pm$ 41.22 & 47.51 $\pm$ 13.08 & 81.25 $\pm$ 32.47 & 0.337 $\pm$ 0.031 & 44.79 $\pm$ 15.57 & 72.06 $\pm$ 7.36 \\
Poisoned client, FedAvg + gate & 10,515 & 88.56 $\pm$ 3.51 & 81.32 $\pm$ 0.69 & 38.82 $\pm$ 2.34 & 0.799 $\pm$ 0.009 & 76.29 $\pm$ 1.14 & 89.56 $\pm$ 0.31 \\
Poisoned client, median & 10,515 & 50.86 $\pm$ 6.05 & 86.54 $\pm$ 0.21 & 18.46 $\pm$ 2.20 & 0.863 $\pm$ 0.003 & 69.16 $\pm$ 13.43 & 98.67 $\pm$ 0.03 \\
Fuzzers on 2 clients & 10,515 & 51.15 $\pm$ 6.93 & 86.46 $\pm$ 0.22 & 18.47 $\pm$ 2.19 & 0.862 $\pm$ 0.003 & 74.36 $\pm$ 0.52 & 98.22 $\pm$ 1.02 \\
Fuzzers on 3 clients & 10,515 & 79.42 $\pm$ 1.40 & 82.72 $\pm$ 0.22 & 33.56 $\pm$ 0.83 & 0.818 $\pm$ 0.003 & 77.03 $\pm$ 0.27 & 98.78 $\pm$ 0.02 \\
5$\times$ local steps & 10,515 & 71.41 $\pm$ 11.17 & 84.15 $\pm$ 1.92 & 28.24 $\pm$ 6.83 & 0.835 $\pm$ 0.022 & 53.93 $\pm$ 1.34 & 97.10 $\pm$ 0.08 \\
3 clients + 5$\times$ steps & 10,515 & 87.19 $\pm$ 4.38 & 67.04 $\pm$ 12.69 & 37.36 $\pm$ 1.44 & 0.665 $\pm$ 0.119 & 54.24 $\pm$ 2.40 & 97.63 $\pm$ 1.09 \\
\bottomrule\end{tabular}\end{table*}

\begin{table}[t]\centering
\caption{Per-flow decisions, ToN\_IoT (real IPs), default cost table (frozen model).}
\label{tab:policy_ton}
\begin{tabular}{lcccc}\toprule
Policy & False block (\%) & Benign alert (\%) & Attack missed (\%) & Cost/flow \\ \midrule
Threshold ($p\geq0.9$) & 1.24 $\pm$ 0.00 & 0.69 $\pm$ 0.00 & 0.92 $\pm$ 0.00 & 0.152 $\pm$ 0.000 \\
Cost-aware (Bayes) & 1.28 $\pm$ 0.00 & 2.40 $\pm$ 0.00 & 0.39 $\pm$ 0.00 & 0.126 $\pm$ 0.000 \\
\bottomrule\end{tabular}\end{table}

\begin{table}[t]\centering
\caption{Four-agent stream simulation, ToN\_IoT (real IPs) (3 seeds). Dropped = benign flows refused by an active block; unblocked = attack flows not blocked or dropped.}
\label{tab:agent_ton}
\begin{tabular}{lccc}\toprule
Agent configuration & Benign dropped (\%) & Attack unblocked (\%) & Rollbacks \\ \midrule
A: threshold rule, permanent block & 49.79 $\pm$ 7.13 & 0.029 $\pm$ 0.008 & 0 $\pm$ 0 \\
B: cost-aware, permanent block & 51.07 $\pm$ 6.75 & 0.028 $\pm$ 0.006 & 0 $\pm$ 0 \\
C: + graded blocks (TTL) & 28.63 $\pm$ 5.67 & 0.083 $\pm$ 0.021 & 0 $\pm$ 0 \\
D: + analyst rollback \& exceptions & 9.40 $\pm$ 0.70 & 0.072 $\pm$ 0.013 & 34 $\pm$ 0 \\
E: + peer intel w/ retraction (full) & 12.50 $\pm$ 1.91 & 0.019 $\pm$ 0.005 & 42 $\pm$ 4 \\
\bottomrule\end{tabular}\end{table}

\begin{table}[t]\centering
\caption{Per-flow decisions, UNSW-NB15 (synthetic sessions), default cost table (frozen model).}
\label{tab:policy_unsw}
\begin{tabular}{lcccc}\toprule
Policy & False block (\%) & Benign alert (\%) & Attack missed (\%) & Cost/flow \\ \midrule
Threshold ($p\geq0.9$) & 1.69 $\pm$ 0.00 & 12.62 $\pm$ 0.00 & 11.59 $\pm$ 0.00 & 1.012 $\pm$ 0.000 \\
Cost-aware (Bayes) & 2.31 $\pm$ 0.00 & 31.53 $\pm$ 0.00 & 2.36 $\pm$ 0.00 & 0.776 $\pm$ 0.000 \\
\bottomrule\end{tabular}\end{table}

\begin{table}[t]\centering
\caption{Four-agent stream simulation, UNSW-NB15 (synthetic sessions) (3 seeds). Dropped = benign flows refused by an active block; unblocked = attack flows not blocked or dropped.}
\label{tab:agent_unsw}
\begin{tabular}{lccc}\toprule
Agent configuration & Benign dropped (\%) & Attack unblocked (\%) & Rollbacks \\ \midrule
A: threshold rule, permanent block & 22.05 $\pm$ 0.60 & 0.381 $\pm$ 0.044 & 0 $\pm$ 0 \\
B: cost-aware, permanent block & 28.35 $\pm$ 0.62 & 0.285 $\pm$ 0.022 & 0 $\pm$ 0 \\
C: + graded blocks (TTL) & 1.42 $\pm$ 0.06 & 1.267 $\pm$ 0.119 & 0 $\pm$ 0 \\
D: + analyst rollback \& exceptions & 1.21 $\pm$ 0.02 & 1.113 $\pm$ 0.160 & 882 $\pm$ 8 \\
E: + peer intel w/ retraction (full) & 2.14 $\pm$ 0.08 & 0.132 $\pm$ 0.032 & 1076 $\pm$ 26 \\
\bottomrule\end{tabular}\end{table}

\begin{table}[t]\centering
\caption{Threshold-free view of federation (official UNSW-NB15 test set, 3 seeds; EC2). AUROC over all flows, Fuzzers-vs-benign AUROC, and Fuzzers recall at a fixed 10\% benign FPR (threshold set on the benign test flows), before (round 0) and after (round 4) federation.}
\label{tab:operating}
\begin{tabular}{llccc}\toprule
Variant & Round & AUROC & Fuzzers AUROC & Fuzzers rec.\ @10\% FPR (\%) \\ \midrule
LoRA $r{=}4$ (default) & 0 & 0.962 $\pm$ 0.000 & 0.823 $\pm$ 0.000 & 37.99 $\pm$ 0.00 \\
 & 4 & 0.959 $\pm$ 0.001 & 0.818 $\pm$ 0.003 & 34.83 $\pm$ 1.40 \\
LoRA $r{=}8$ & 0 & 0.962 $\pm$ 0.000 & 0.823 $\pm$ 0.000 & 37.99 $\pm$ 0.00 \\
 & 4 & 0.958 $\pm$ 0.001 & 0.818 $\pm$ 0.003 & 34.40 $\pm$ 1.45 \\
Full-model FedAvg & 0 & 0.962 $\pm$ 0.000 & 0.823 $\pm$ 0.000 & 37.99 $\pm$ 0.00 \\
 & 4 & 0.950 $\pm$ 0.003 & 0.809 $\pm$ 0.003 & 30.19 $\pm$ 0.51 \\
Median aggregation & 0 & 0.962 $\pm$ 0.000 & 0.823 $\pm$ 0.000 & 37.99 $\pm$ 0.00 \\
 & 4 & 0.959 $\pm$ 0.001 & 0.817 $\pm$ 0.001 & 34.63 $\pm$ 0.98 \\
Fuzzers on 2 clients & 0 & 0.962 $\pm$ 0.000 & 0.823 $\pm$ 0.000 & 37.99 $\pm$ 0.00 \\
 & 4 & 0.959 $\pm$ 0.001 & 0.818 $\pm$ 0.001 & 34.69 $\pm$ 0.60 \\
Fuzzers on 3 clients & 0 & 0.962 $\pm$ 0.000 & 0.823 $\pm$ 0.000 & 37.99 $\pm$ 0.00 \\
 & 4 & 0.957 $\pm$ 0.000 & 0.818 $\pm$ 0.001 & 34.48 $\pm$ 0.31 \\
5$\times$ local steps & 0 & 0.962 $\pm$ 0.000 & 0.823 $\pm$ 0.000 & 37.99 $\pm$ 0.00 \\
 & 4 & 0.956 $\pm$ 0.002 & 0.817 $\pm$ 0.001 & 33.77 $\pm$ 0.86 \\
3 clients + 5$\times$ steps & 0 & 0.962 $\pm$ 0.000 & 0.823 $\pm$ 0.000 & 37.99 $\pm$ 0.00 \\
 & 4 & 0.826 $\pm$ 0.115 & 0.820 $\pm$ 0.006 & 36.12 $\pm$ 2.98 \\
\bottomrule\end{tabular}\end{table}

\begin{table}[t]\centering
\caption{Federated Fuzzers recall vs.\ the frozen backbone with only its decision threshold moved to the same benign FPR (official UNSW-NB15 test set; EC2). The frozen model needs no federation at all.}
\label{tab:threshold}
\begin{tabular}{lccc}\toprule
Federated model & FPR (\%) & Fuzzers rec.\ (\%) & Frozen, threshold moved (\%) \\ \midrule
LoRA $r{=}4$, 4 rounds & 18.28 $\pm$ 0.83 & 50.48 $\pm$ 3.16 & 51.62 \\
LoRA $r{=}4$, Lambda path (4 seeds) & 19.21 $\pm$ 0.61 & 54.50 $\pm$ 1.92 & 53.37 \\
LoRA $r{=}8$, 4 rounds & 21.24 $\pm$ 2.66 & 59.12 $\pm$ 6.45 & 56.96 \\
Full-model FedAvg, 4 rounds & 33.08 $\pm$ 1.80 & 78.04 $\pm$ 3.34 & 84.00 \\
Fuzzers on 2 clients, 4 rounds & 18.47 $\pm$ 2.19 & 51.15 $\pm$ 6.93 & 51.93 \\
5$\times$ local steps, 4 rounds & 28.24 $\pm$ 6.83 & 71.41 $\pm$ 11.17 & 71.68 \\
Fuzzers on 3 clients, 4 rounds & 33.56 $\pm$ 0.83 & 79.42 $\pm$ 1.40 & 85.09 \\
3 clients + 5$\times$ steps, 4 rounds & 37.36 $\pm$ 1.44 & 87.19 $\pm$ 4.38 & 89.69 \\
\bottomrule\end{tabular}\end{table}

\begin{table*}[t]\centering
\caption{Sensitivity of the full agent to the analyst (3 seeds; EC2). The analyst reviews a random 50\% of blocks and gives the wrong verdict with the stated probability; a wrong verdict on an attacker rolls its block back and grants it the learned exception.}
\label{tab:analyst}
\begin{tabular}{lcccc}\toprule
 & \multicolumn{2}{c}{ToN\_IoT} & \multicolumn{2}{c}{UNSW-NB15} \\
Configuration & Benign dropped (\%) & Attack unblocked (\%) & Benign dropped (\%) & Attack unblocked (\%) \\ \midrule
C: no analyst & 28.63 $\pm$ 5.67 & 0.083 $\pm$ 0.021 & 1.42 $\pm$ 0.06 & 1.267 $\pm$ 0.119 \\
E: oracle analyst, 30-min exception & 12.50 $\pm$ 1.91 & 0.019 $\pm$ 0.005 & 2.14 $\pm$ 0.08 & 0.132 $\pm$ 0.032 \\
E: 50\% reviewed, 10\% wrong & 23.97 $\pm$ 4.54 & 4.543 $\pm$ 1.583 & 4.05 $\pm$ 0.32 & 3.067 $\pm$ 1.217 \\
E: 50\% reviewed, 20\% wrong & 31.77 $\pm$ 8.34 & 11.785 $\pm$ 3.186 & 4.13 $\pm$ 0.27 & 12.362 $\pm$ 1.731 \\
E: oracle, 5-min exception & 18.56 $\pm$ 1.58 & 0.018 $\pm$ 0.006 & 2.38 $\pm$ 0.10 & 0.132 $\pm$ 0.032 \\
E: 10\% wrong, 5-min exception & 29.16 $\pm$ 4.45 & 1.474 $\pm$ 0.228 & 4.25 $\pm$ 0.14 & 1.063 $\pm$ 0.396 \\
E: 20\% wrong, 5-min exception & 34.46 $\pm$ 6.44 & 4.893 $\pm$ 0.235 & 4.40 $\pm$ 0.36 & 3.049 $\pm$ 0.783 \\
\bottomrule\end{tabular}\end{table*}

\begin{table*}[t]\centering
\caption{Reference models on the same test sets and eight features (3 seeds each; trained and evaluated on EC2). Neither is federated; they bound what the frozen backbone and the federated system are compared against.}
\label{tab:refs}
\begin{tabular}{lcccc}\toprule
Model (training data) & UNSW acc. (\%) & UNSW FPR (\%) & Fuzzers rec. (\%) & ToN acc. (\%) \\ \midrule
Gradient-boosted trees (pretraining pool, Fuzzers excluded) & 90.10 $\pm$ 0.17 & 8.92 $\pm$ 0.40 & 34.27 $\pm$ 0.25 & 99.55 $\pm$ 0.01 \\
Centralized Transformer (all device data incl.\ Fuzzers) & 82.39 $\pm$ 0.45 & 36.38 $\pm$ 1.65 & 89.76 $\pm$ 1.74 & 98.88 $\pm$ 0.07 \\
Backbone v1, frozen (pretraining pool) & 87.19 & 14.31 & 41.06 & 98.84 \\
Backbone v1 + 4 LoRA rounds (federated) & 86.46 $\pm$ 0.07 & 19.21 $\pm$ 0.61 & 54.50 $\pm$ 1.92 & -- \\
\bottomrule\end{tabular}\end{table*}
 (generated by
% paper_figures/build_results.py from the EC2 logs -- do not hand-edit).
% Requires: \usepackage{amsmath,booktabs,graphicx}
% =====================================================================

\section{System Model and Methodology}

\subsection{Threat model and setting}
$K{=}4$ edge nodes (AWS EC2 t3.small, 2 vCPU, 2\,GB RAM) each observe network
flows and hold a private, non-IID labelled partition (Dirichlet $\alpha{=}0.5$
over attack categories). Devices 1--2 hold UNSW-NB15 flows and devices 3--4
hold ToN\_IoT flows. One UNSW-NB15 category (Fuzzers) is withheld from central pretraining and
from every device except device~1, emulating an attack that only one site has
labelled examples of. Raw flows never leave a device; only adapter weights and
incident metadata cross the network. Coordination runs on AWS serverless
services (API Gateway, Lambda, DynamoDB, S3); nodes authenticate with
per-device API keys. Fig.~\ref{fig:arch} shows the architecture.

\subsection{Flow representation and backbone}
Each flow is mapped to a cross-dataset canonical schema of five continuous
features (duration, source/destination bytes, source/destination packets)
and three categorical features (protocol, service, connection state).
Continuous features are transformed as $\tilde x_j = (\log(1+x_j)-\mu_j)/\sigma_j$
with $\mu_j,\sigma_j$ frozen centrally so all devices share one encoding.
Backbone v2 embeds every continuous feature as its own token,
$\mathbf t_j = \tilde x_j \mathbf w_j + \mathbf b_j$ ($\mathbf w_j,\mathbf b_j\in\mathbb R^{d}$),
and every categorical feature through an embedding table, giving
$n{=}8$ tokens of width $d{=}32$. Backbone v1 instead projects all five
continuous features into a single token ($n{=}4$). A two-layer Transformer encoder
(4 heads, FFN width $4d$) is followed by mean pooling and a linear head.
Unless stated otherwise, all federated, ablation and agent experiments use
backbone v1; v2 is evaluated as a backbone and in one federated
replication (Table~\ref{tab:fed_rounds_v2}).

\subsection{Parameter-efficient federated adaptation}
After central pretraining the backbone is frozen. Each FFN projection
$W_0\in\mathbb R^{m\times n}$ receives a low-rank update
\begin{equation}
W = W_0 + \frac{\alpha}{r} B A,\quad A\in\mathbb R^{r\times n},\; B\in\mathbb R^{m\times r},
\end{equation}
with $r{=}4$, $\alpha{=}8$, and $B$ initialised to zero. Only $\{A,B\}$ and
the classification head are trained and communicated
($\approx$2.6K parameters, 10.3\,KB of float32 per upload).
The server aggregates with FedAvg,
\begin{equation}
\theta^{t+1} = \sum_{k\in S_t} \frac{n_k}{\sum_{j\in S_t} n_j}\,\theta_k^{t},
\end{equation}
where $S_t$ is the set of devices that reported in round $t$. A round closes
when all $K$ devices report or, via a scheduled straggler handler, after a
timeout with at least one participant. FedProx~\cite{li2020fedprox} is evaluated as an
alternative local objective,
$\ell_k(\theta) + \frac{\mu}{2}\lVert\theta-\theta^{t}\rVert^2$, $\mu{=}0.01$.

\subsection{Local training}
Each local update runs a fixed budget of $E{\cdot}S = 5{\times}20$ minibatch
steps (batch 64), independent of buffer size. Minibatches are sampled with
probability proportional to a clipped per-example loss,
$w_i = \min(\ell_i, 3) + 0.3$, recomputed each epoch (online hard-example
mining). Each step minimises clean plus FGSM-perturbed cross-entropy,
\begin{equation}
\begin{aligned}
\mathcal L ={}& \mathrm{CE}(f(x),y)\\
&+ \mathrm{CE}\big(f(x+\epsilon\,\mathrm{sign}(\nabla_x \mathrm{CE}(f(x),y))),y\big),
\end{aligned}
\end{equation}
with $\epsilon{=}0.05$, applied to the normalised continuous features. A shuffled 15\% of each
local buffer is held out for the device's own validation.

\subsection{The edge agent}
Each node runs a perceive--decide--act--reflect loop. Unlike a fixed pipeline,
the agent chooses between actions under uncertainty, verifies what it
receives, and adapts from feedback.

\textbf{Cost-aware action selection.} For a flow with
$p = P(\text{attack}\mid x)$ the agent takes
\begin{equation}
\begin{aligned}
a^\ast = \arg\min_{a\in\{\textsc{noop},\textsc{alert},\textsc{block}\}}
\;&(1-p)\,C(a,\text{benign})\\
&+ p\,C(a,\text{attack}),
\end{aligned}
\end{equation}
with default costs $C(\textsc{noop})=(0,10)$, $C(\textsc{alert})=(1,3)$,
$C(\textsc{block})=(20,0)$ (benign, attack). These are deployment
assumptions, not learned values, and must be set per site. The
original rule ($\textsc{block}$ iff $p\ge0.9$) is kept as a baseline.

\textbf{Graded response.} A block on source IP $s$ raises its level
$\lambda_s$. Levels $\lambda_s<4$ are temporary with TTL $60\cdot2^{\lambda_s-1}$\,s,
and level 4 is permanent. Expired blocks lift automatically.

\textbf{Analyst feedback.} A block an analyst marks as a false positive is
rolled back. The host gets a 30-minute exception (block verdicts downgraded
to alerts), and the benign-side block cost is updated multiplicatively,
$C(\textsc{block},\text{benign}) \leftarrow \min(1.1\,C, 100)$; a confirmed
block lowers it by 5\% (floor 0.1).

\textbf{Shared threat intelligence with retraction.} Every 15\,s an agent
pulls recent $\textsc{block}$/$\textsc{rollback}$ records of its peers from
the cloud incident store. An IP whose latest record is a peer's block is
escalated faster: alerts become blocks, starting one level higher. A
rollback by any device retracts the IP everywhere.

\textbf{Verified model adoption.} When a new global model $\theta^{t+1}$
arrives, the agent evaluates it and its current model on its local held-out
set and adopts it only if
\begin{equation}
\begin{aligned}
&\mathrm{Acc}(\theta^{t+1}) \ge \mathrm{Acc}(\theta_k) - 0.01 \;\wedge\\
&\mathrm{FPR}(\theta^{t+1}) \le \mathrm{FPR}(\theta_k) + 0.02 .
\end{aligned}
\end{equation}
Otherwise it keeps its model and logs a rejection.

\textbf{Autonomous retraining.} A device retrains when the rolling mean
of its top-class confidence over the last 50 flows falls below 0.6
(minimum 20 flows), or when its accuracy on its own labelled buffer falls
below 0.90. Self-evaluation runs at most once per minute, and retraining at
most once per 5 minutes.

\textbf{Explanation.} Every \textsc{block}/\textsc{alert} carries an
attention-rollout attribution~\cite{abnar2020},
$\hat A^{(l)} = 0.5A^{(l)}+0.5I$ (row-normalised),
$R=\hat A^{(L)}\cdots\hat A^{(1)}$, with token importance given by the column
sums of $R$. This is a heuristic attribution, not a causal explanation.
An optional LLM triage worker (Claude on Amazon Bedrock) summarises
incidents for analysts asynchronously, off the detection path. It receives
only incident metadata and is not evaluated in this paper.

\section{Experimental Setup}
All reported numbers were produced on the four EC2 t3.small nodes
(ap-south-1). The federated round experiments execute the deployed Lambda
handler code on the nodes, with DynamoDB and S3 replaced by moto mocks, so that
rounds are reproducible and seeded. The ablation runs execute the same FedAvg
arithmetic in-process. The live backend (API Gateway, Lambda, DynamoDB, S3,
per-device API keys) was exercised separately by the deployed agents: between
21 September and 16 October 2026 it closed 12{,}610 rounds. Only 22 of them
had all four devices; the scheduled straggler handler closed the rest with one
to three participants. This validates the coordination path under real
partial participation but is not used for any accuracy number. Each
configuration is repeated over 3--4 seeds and reported as mean $\pm$ sample
standard deviation. Test sets are never used for training or model selection:
the official UNSW-NB15 testing set (82{,}332 flows; 6{,}062 Fuzzers) and a
stratified 15\% ToN\_IoT split (31{,}656 flows). Backbone checkpoints are
selected on a random 15\% split of the pretraining pool. Besides
default-threshold metrics we report threshold-free ones: AUROC, Fuzzers-vs-benign
AUROC, and Fuzzers recall at a fixed benign FPR. The threshold for the last is
the corresponding quantile of the model's scores on the benign test flows; it
is an evaluation device and is never used by the deployed system. Reference
models (gradient-boosted trees, scikit-learn \texttt{HistGradientBoosting}, and
a centrally trained Transformer) use the same eight features and test sets.

\paragraph{Agent simulation} The real agent, actuator and policy code are
driven on a simulated clock, with one flow every 0.05\,s distributed
round-robin over four agents. ToN\_IoT keeps its real source IPs (47 hosts).
UNSW-NB15's IPs are synthetic, so flows are grouped into synthetic sessions
(300 benign hosts, 5 attacker IPs per category); this grouping is an
assumption. The analyst reviews blocks after 300\,s. In the main runs it
reviews every block and is always correct, an upper bound on what feedback can
achieve. A sensitivity study then has it review a random 50\% of blocks with a
10\% or 20\% chance of a wrong verdict, and varies the learned-exception
duration (30 vs.\ 5 minutes).

\section{Limitations}
\begin{itemize}
\item On the official UNSW-NB15 split the backbone reaches 87.7\% accuracy,
below the $>$96\% often reported on random splits (97.9\% on our own random
validation split). Gradient-boosted trees on the same eight features reach
90.1\%, so about 2.4 points of the gap are attributable to the Transformer
itself; the rest is distribution shift between the official files.
\item Federated adaptation did not improve class separation in our setting:
its recall gain equals a threshold shift of the frozen model
(Table~\ref{tab:threshold}), and UNSW FPR rises with every round. This held
with one to three of four clients holding the new class and with 100 or 500
local steps per round, over four rounds. More clients, more rounds or other
local objectives may differ.
\item FedAvg weights clients by self-reported $n_k$. A malicious client can
inflate its weight. The adoption gate limits but does not remove this.
\item Device partitions are drawn from the same pool as pretraining, except
device~1's Fuzzers rows. The (label-free) normalisation statistics were computed
before the ToN\_IoT test split was carved out. Recomputing them without the test
rows shifts each feature mean by at most $0.058\sigma$ and each standard deviation
by at most 2.4\%; the official UNSW-NB15 test file was never included.
\item The headline agent results assume an oracle analyst; with an
imperfect analyst, attack leakage rises sharply (Table~\ref{tab:analyst}).
UNSW-NB15 agent results use synthetic host sessions. Live capture was
validated only with loopback scans.
\item Four clients only. Round experiments use a mocked cloud backend for
reproducibility; the live backend was validated operationally, not for accuracy.
\end{itemize}
